% generated by regen_study_tables.py -- do not edit
% source: 2026-09-25-r2d2-prior/tables/tier_d_calibration.csv, 2026-09-25-r2d2-prior/REPORT.md section G0 decision
\begin{table}[tbp]
\centering
\caption{Calibration of the R2-D2 prior at stage 0 (directory \SOcode{2026-09-25-r2d2-prior}, file \SOcode{tier\_d\_calibration.csv}; scope: the prior alone at dimension 100, no optimization runs; the profile and trend rows of the file are not shown). Each point's quartiles of the prior's active-coordinate count are compared with the half-Cauchy cells'; every point matches the target median, and \SOcode{C\_tied}, median-matched at its fixed b = 0.5, matches the median only (quartiles 30/46 against 17/64); alpha is k times the dimension, and the residual is the larger of the two outer-quartile gaps to the target. Gate G0 chose the reference (untied) form at \SOcode{C\_untied}, (a, b, k) = (1.577790731256835, 0.8043916352200708, 0.49942145555129697): its quartiles 14/37/61 against the target's 17/37/64, residuals 3 and 3 on the lower and upper quartile (continuous 2.93). The study tables beside this one were read from snapshot \SOcode{2026-10-02-1420}.}\label{app:so:study:tab:calibration}
\scriptsize\setlength{\tabcolsep}{2pt}
\begin{tabular}{lrrrrrrrrr}
\toprule
point & a & b & k & alpha & \shortstack{lower\\ quartile} & median & \shortstack{upper\\ quartile} & residual & \shortstack{R2 mass\\ below 0.6} \\
\midrule
half-Cauchy target & --- & --- & --- & --- & 17 & 37 & 64 & 0 & --- \\
\SOcode{C\_untied} (chosen) & 1.5778 & 0.8044 & 0.4994 & 49.94 & 14 & 37 & 61 & 3 & 0.3705 \\
\SOcode{C\_untied\_k1} & 3.0487 & 1.6738 & 0.9981 & 99.81 & 14 & 37 & 62 & 3 & 0.3866 \\
\SOcode{C\_tied} & 9.3786 & 0.5000 & 0.0938 & 9.38 & 30 & 37 & 46 & 18 & 0.0022 \\
\bottomrule
\end{tabular}
\end{table}
